\documentclass[11pt]{article}
\usepackage[a4paper,margin=2.05cm]{geometry}
\usepackage[T1]{fontenc}
\usepackage{lmodern}
\usepackage{amsmath,amssymb}
\usepackage{xcolor}
\usepackage{microtype}
\usepackage[hidelinks]{hyperref}
\setlength{\emergencystretch}{3em}

\definecolor{ink}{HTML}{172033}
\definecolor{accent}{HTML}{385B7A}
\definecolor{soft}{HTML}{EEF5FA}
\color{ink}

\title{Arbitrary-dimensional CPWL density in Hessian--Schatten energy}
\author{Literature answer to Conjecture 1 in arXiv:2210.04077}
\date{17 August 2026}

\begin{document}
\maketitle

\begin{center}
\fcolorbox{accent}{soft}{\parbox{0.91\textwidth}{
\textbf{Status: literature-already-answered, full intended conjecture.}
For every $d$, continuous piecewise-linear functions are dense, in
$L^1((0,1)^d)$ and in Schatten-$1$ Hessian energy, among functions of bounded
Hessian--Schatten variation.  Ambrosio--Brena--Conti explicitly identify their
Theorem~2.4 as a positive answer to the source's Conjecture~1.}}
\end{center}

\section*{Source conjecture}
Ambrosio--Aziznejad--Brena--Unser prove in Theorem~21 that, for
$\Omega=(0,1)^2$, every function $f$ of bounded Hessian--Schatten variation
has CPWL approximants $f_k$ satisfying
\[
 f_k\longrightarrow f \quad\hbox{in }L^\infty(\Omega),
 \qquad
 |D_1^2 f_k|(\Omega)\longrightarrow |D_1^2 f|(\Omega).
\]
Their Conjecture~1, on source PDF page~14, asks whether the density result
extends to every hypercube $(0,1)^n$ \cite{AABU}.  The arXiv v1 wording says
only that Theorem~21 should remain valid in arbitrary dimension.  The
published version makes the intended formulation precise: in dimensions
$n\geq3$, the $L^\infty$ topology is replaced by $L^1$ (and it notes
$L^{n/(n-2)}$ as another natural ambient topology).  Its conjecture footnote
also records Sergio Conti's announced proof.

\section*{Explicit later solution}
Ambrosio--Brena--Conti state immediately before Theorem~2.4, on supporting PDF
page~14, that the theorem ``gives a positive answer'' to the source's
Conjecture~1 \cite{ABC}.  Their precise result is:

\medskip
\noindent
For every $d$ and $\Omega=(0,1)^d$, if $f\in L^1_{\mathrm{loc}}(\Omega)$ has
bounded Hessian--Schatten variation, then there are
$f_k\in\mathrm{CPWL}(\Omega)$ such that
\[
 f_k\longrightarrow f \quad\hbox{in }L^1(\Omega),
 \qquad
 |D_1^2 f_k|(\Omega)\longrightarrow |D_1^2 f|(\Omega).
\]

This is exactly the published arbitrary-dimensional conjecture.  The
supporting authors knew the identification: they cite Conjecture~1 and
Theorem~21 of the source explicitly.  This is therefore a retrieved literature
answer, not an agent-generated theorem.

\section*{Mechanism and scope}
The supporting paper first proves its global Theorem~2.2 for smooth targets in
arbitrary dimension.  It constructs triangulations whose local orientations
follow the eigendirections of the target Hessian and controls the energy cost
of gluing differently oriented local meshes.  An extension lemma,
mollification, and a diagonal argument then yield the localized Theorem~2.4
for arbitrary finite-energy targets.

The answer is full for the intended $L^1$ formulation.  It should not be read
as an $L^\infty$ density theorem in dimensions $d\geq3$.  In dimension two,
the source's $L^\infty$ result remains stronger in its function topology.

\section*{Evidence checked}
The identification was checked against the local arXiv sources and rebuilt
PDFs for both papers, as well as the publishers' official article text.  The
decisive locations are Conjecture~1 on source PDF page~14 and Theorem~2.4,
including its explicit positive-answer sentence, on supporting PDF page~14.

{\raggedright
\begin{thebibliography}{9}
\bibitem{AABU}
L. Ambrosio, S. Aziznejad, C. Brena, and M. Unser,
\emph{Linear Inverse Problems with Hessian--Schatten Total Variation},
arXiv:2210.04077; Calc. Var. Partial Differential Equations \textbf{63},
article 9 (2024),
\href{https://doi.org/10.1007/s00526-023-02611-6}{doi:10.1007/s00526-023-02611-6}.

\bibitem{ABC}
L. Ambrosio, C. Brena, and S. Conti,
\emph{Functions with Bounded Hessian--Schatten Variation: Density,
Variational, and Extremality Properties},
arXiv:2302.12554; Arch. Ration. Mech. Anal. \textbf{247}, article 111 (2023),
\href{https://doi.org/10.1007/s00205-023-01938-w}{doi:10.1007/s00205-023-01938-w}.
\end{thebibliography}
\par}

\end{document}
